\section*{Chapter \ref{ch:qm:high-frequency-econometrics} --- High-Frequency Econometrics}

\begin{solution}{exo:qm:high-frequency-econometrics:1}
$\E[\mathrm{RV}] = 10^{-4} + 2 \times 23\,400 \times (2 \times 10^{-4})^2 = 10^{-4} + 1.87 \times 10^{-3} = 1.97 \times 10^{-3}$: a daily volatility of 4.4\% for a true 1\%.
\end{solution}

\begin{solution}{exo:qm:high-frequency-econometrics:2}
$2\sqrt{0.0004} = 0.04$ dollars, four cents.
\end{solution}

\begin{solution}{exo:qm:high-frequency-econometrics:3}
The mid does not jump between bid and ask with the trade sign, so the bounce is gone; but it moves only when the quotes cross a grid point, so it is the efficient price
rounded to half-ticks, and the rounding error is noise.
\end{solution}

\begin{solution}{exo:qm:high-frequency-econometrics:4}
$\mathrm{IQ} = (10^{-4})^2 = 10^{-8}$ and $\omega^4 = (4 \times 10^{-8})^2 = 1.6 \times 10^{-15}$, so $n^* = (10^{-8}/6.4 \times 10^{-15})^{1/3} = 116$ returns, one every 202 seconds.
\end{solution}

\begin{solution}{exo:qm:high-frequency-econometrics:5}
An observed return is $r_i + u_i - u_{i-1}$ with variance $\sigma^2\Delta + 2\omega^2$; two consecutive ones share $u_{i-1}$ with opposite signs, so their covariance is
$-\omega^2$. For the chapter's stock at one second, $\sigma^2\Delta = 2.48 \times 10^{-4}/23\,400 = 1.06 \times 10^{-8}$ and $\omega^2 = 2.56 \times 10^{-8}$: $-2.56/(1.06 + 5.12) =
-0.41$, as measured in the simulation ($-0.41$).
\end{solution}

\begin{solution}{exo:qm:high-frequency-econometrics:6}
With independent trading, $0.2 \times 1/15 = 1.3\%$ of one-second intervals contain a trade in both; in the others, at least one previous-tick return is zero, which is why the
one-second correlation collapses.
\end{solution}

\begin{solution}{exo:qm:high-frequency-econometrics:7}
\texttt{signature()} and \texttt{mse\_by\_interval()}: trades 60.1\%, 35.1\%, 26.1\%, 25.4\% at 1, 5, 60, 300 seconds; mids 36.8\% at one second; efficient 25.0\% everywhere. The
root-mean-squared relative error is 4.80 at one second and smallest near one minute (0.11 at 60 seconds, 0.12 at 120).
\end{solution}

\begin{solution}{exo:qm:high-frequency-econometrics:8}
One-second returns of two stocks rarely contain trades in both, so previous-tick realised correlation is biased toward zero (the \omterm{def:qm:high-frequency-econometrics:epps}{Epps effect}): 0.12 is an artefact of
the sampling, not a property of the pair. Use Hayashi--Yoshida, or sample at several minutes, or a synchronised noise-robust \omterm{def:qm:estimation:estimator}{estimator}, and compare.
\end{solution}

\begin{solution}{pb:qm:high-frequency-econometrics:1}
\textbf{1.} 60.1\%, 35.1\%, 26.1\% and 25.4\%.
\textbf{2.} 36.8\% from mid quotes at one second; 25.0\% from the efficient price at every interval.
\textbf{3.} Estimated from $\mathrm{RV}/2n$: 1.75 basis points; true: 1.60 (a trade is uniformly within one tick of the efficient price, a tick being 2.8 basis points).
\textbf{4.} 1.14 cents, against a true one-cent spread.
\textbf{5.} Each one-second return carries a noise difference of variance $2\omega^2$, about five times the efficient variance per second, and 23\,400 of them add up: RV is 5.8 times
the \omterm{def:qm:volatility-models:rv}{integrated variance}.
\textbf{6.} $\omega^2/\mathrm{IV} = 1.24 \times 10^{-4}$.
\textbf{7.} $n^* = 254$, one return every 92 seconds (82 with the true noise).
\textbf{8.} 4.80 at one second, 0.98 at five, 0.18 at thirty, 0.11 at sixty, 0.12 at 120, 0.14 at 300 and 0.34 at 1\,800 seconds.
\textbf{9.} It uses 254 of 23\,400 returns, about 1\%.
\textbf{10.} $n^* \propto \mathrm{IQ}^{1/3}$ and IQ scales with the square of the variance, so doubling the variance multiplies $n^*$ by $2^{2/3} = 1.59$: sample about every 58 seconds.
\textbf{11.} 25.2\%, 25.1\% and 25.1\%, with daily relative errors of 11.9\%, 7.2\% and 4.8\% (five-minute RV: 25.4\% and 14.3\%).
\textbf{12.} Realised correlation 0.03 at one second, 0.21 at ten, 0.49 at sixty, 0.59 at 300 seconds, for a true 0.6.
\textbf{13.} Hayashi--Yoshida 0.594 (standard deviation 0.030 across days); \omterm{def:qm:high-frequency-econometrics:hy}{refresh-time sampling} 0.48.
\textbf{14.} At five minutes RV rises to 1.43 times the \omterm{def:qm:volatility-models:rv}{integrated variance} and \omterm{def:qm:high-frequency-econometrics:bipower}{bipower variation} to 1.16 (from 1.00 and 1.00), for a jump worth 29\% of the day's \omterm{def:qm:brownian-motion:qv}{quadratic
variation}.
\textbf{15.} \omterm{def:qm:high-frequency-econometrics:bipower}{Bipower variation} needs fine sampling to isolate the jump, and fine sampling brings the noise back; noise-robust \omterm{def:qm:estimation:estimator}{estimators}, in turn, count jumps as variance.
\textbf{16.} A noise-robust estimate (\omterm{def:qm:high-frequency-econometrics:kernel}{realised kernel} or pre-averaging) on the latest window, blended with a daily forecast (chapter~18).
\textbf{17.} Hayashi--Yoshida or a noise-robust synchronised \omterm{def:qm:estimation:estimator}{estimator}, never previous-tick one-second returns.
\textbf{18.} Noise that depends on the trade sign and the spread, clustered trading, and time-varying volatility within the day.
\textbf{19.} Named result: \emph{the volatility that depended on the clock}: with a noise-to-signal ratio of $1.24 \times 10^{-4}$, the optimal sampling interval is about 90 seconds (92
from the estimated noise, 82 from the true), and it avoids a 5.8-fold overstatement of the variance by one-second \omterm{def:qm:volatility-models:rv}{realised variance} (60\% volatility for a true 25\%).
\textbf{20.} Because every extra observation adds the same noise variance while the efficient variation it adds shrinks, so beyond some frequency the sum measures the noise.
\end{solution}

\begin{solution}{iq:qm:high-frequency-econometrics:1}
\omterm{def:qm:high-frequency-econometrics:noise}{Microstructure noise}: \omterm{def:qm:high-frequency-econometrics:noise}{bid--ask bounce} and price discreteness add a noise difference to every return, whose variance does not shrink with the interval, so \omterm{def:qm:volatility-models:rv}{realised variance}
from tick data counts $2n$ noise variances. Five-minute sampling makes $n$ small enough that the noise is negligible.
\iqlookfor{The $2n\omega^2$ bias and the \omterm{def:qm:high-frequency-econometrics:signature}{signature plot}.}
\end{solution}

\begin{solution}{iq:qm:high-frequency-econometrics:2}
\omterm{def:qm:high-frequency-econometrics:roll}{Roll's estimator}: under bounce, consecutive trade-price changes are negatively correlated with covariance $-c^2/4$, so the effective spread is $2\sqrt{-\gamma_1}$. With quotes and trade
signs, compare trade prices with the prevailing mid instead.
\iqlookfor{The bounce model and its assumptions.}
\end{solution}

\begin{solution}{iq:qm:high-frequency-econometrics:3}
The interval that balances the discretisation variance ($2\,\mathrm{IQ}/n$) against the noise bias ($2n\omega^2$): $n^* = (\mathrm{IQ}/4\omega^4)^{1/3}$. It depends on the noise-to-signal
ratio: more volatile days and larger stocks (finer relative grid) allow faster sampling.
\iqlookfor{The trade-off and the cube-root formula.}
\end{solution}

\begin{solution}{iq:qm:high-frequency-econometrics:4}
The \omterm{def:qm:high-frequency-econometrics:epps}{Epps effect}: the two stocks do not trade at the same instants, so on a one-second grid most returns of one are matched with zero returns of the other. Use Hayashi--Yoshida,
\omterm{def:qm:high-frequency-econometrics:hy}{refresh-time sampling} with a noise-robust \omterm{def:qm:estimation:estimator}{estimator}, or coarser sampling.
\iqlookfor{Asynchronicity as the cause, and a covariance \omterm{def:qm:estimation:estimator}{estimator} that does not interpolate.}
\end{solution}

\begin{solution}{iq:qm:high-frequency-econometrics:5}
Compare \omterm{def:qm:volatility-models:rv}{realised variance}, which includes squared jumps, with bipower (or threshold) variation, which estimates only the diffusive part; their difference, standardised,
is a jump test. Use moderate sampling or noise-robust versions, since the noise inflates both.
\iqlookfor{RV versus bipower and the noise caveat.}
\end{solution}

\begin{solution}{iq:qm:high-frequency-econometrics:6}
iid noise makes observed returns an MA(1): their variance is inflated by $2\omega^2$ and their first autocovariance is $-\omega^2$. Adding twice the weighted autocovariances to the variance
cancels the inflation, exactly as a long-run variance sums an MA's autocovariances; smooth weights such as Parzen's keep the estimate positive and handle dependent noise.
\iqlookfor{The MA(1) structure and the HAC analogy.}
\end{solution}
